\emph{(1) Anti-automorphism and the rotor.} Grade involution $\alpha$ is an automorphism and reversal is an anti-automorphism of every Clifford algebra, and they commute (\claimref{std-algebra.involution-properties}), so $\alpha \circ \widetilde{\ }$ satisfies $\alpha(\widetilde{xy}) = \alpha(\widetilde y)\,\alpha(\widetilde x)$; it is not multiplicative, since on $x = e_0$, $y = e_1$ it gives $\alpha(\widetilde{e_0e_1}) = -e_0e_1 \neq e_0e_1 = \alpha(e_0)\alpha(e_1)$. On a vector, reversal is the identity and grade involution is $-1$. The blades $B_1 = e_0e_1$, $B_2 = e_2e_3$, $B_3 = e_4e_5$ commute pairwise and $B_k^2 = -e_{2k-2}^2 e_{2k-1}^2 = -1$ in each case (the products of squares are $(+1)(+1)$, $(-1)(-1)$, $(-1)(-1)$), so $\exp(\theta B_k) = \cos\theta + B_k \sin\theta$ (\claimref{std-algebra.bivector-exponential}) and $\exp(\pi B_k/2) = B_k$; hence $I_6 = B_1B_2B_3 = \exp(\tfrac{\pi}{2}(B_1+B_2+B_3))$ is joined to $1$ by the path of rotors $\exp(\tfrac{t}{2}(B_1+B_2+B_3))$, $t\in[0,\pi]$, and $I_6\widetilde{I_6} = 1$. In even dimension $I_6$ anticommutes with every vector (\claimref{std-algebra.pseudoscalar-commutation}), so $\mathrm{Ad}(I_6)$ is $-1$ on vectors and equals $\alpha$ on the whole algebra (both are automorphisms agreeing on generators). In $\Cl(1,3)$ with $e_0^2 = 1$, $e_{1,2,3}^2 = -1$, $\widetilde{I_4} = I_4$ and $I_4^2 = -1$, so $I_4\widetilde{I_4} = -1$ and $I_4$ is not a rotor; this is the algebraic form of Theorem \claimref{singularity.euclidean-rotor-where-it-lies}(3), that the Euclidean rotation is realised by no real rotor of $\Cl(1,3)$. Checked exactly: anti-multiplicativity on fixed elements, failure of multiplicativity, $B_k^2 = -1$, $I_6\widetilde{I_6} = 1$, $\mathrm{Ad}(I_6) = \alpha$ on generators and on products, $I_4\widetilde{I_4} = -1$.

\emph{(2) The conjugate flag.} If $W$ is maximal isotropic with $W \cap \overline W = 0$ then $\C^6 = W \oplus \overline W$ by dimension; define $J$ as multiplication by $I_7$ on $W$ and by $-I_7$ on $\overline W$ ($I_7$ the complex unit of the formal level). $J$ commutes with $\kappa$, so it is real, $J^2 = -1$, and it is orthogonal: for $x, y \in W$, $B(Jx,Jy) = -B(x,y) = 0$ by isotropy, likewise on $\overline W$, and for $x \in W$, $y \in \overline W$, $B(Jx,Jy) = (I_7)(-I_7)B(x,y) = B(x,y)$. Conversely the $+I_7$ eigenspace of any orthogonal complex structure is maximal isotropic and meets its conjugate in $0$. For the structure $J: e_0 \mapsto e_1$, $e_2 \mapsto e_3$, $e_4 \mapsto e_5$ the eigenspace is spanned by $w_k = e_{2k-2} - I_7e_{2k-1}$, with $Q(w_k) = 0$, $w_aw_b + w_bw_a = 0$, and $w_a\overline{w_b} + \overline{w_b}w_a = \mathrm{diag}(4,-4,-4)$, nondegenerate, so $(W, \overline W)$ is a Witt pair. The null blade $N = w_1w_2w_3$ is annihilated on the left by each $w_k$ and is the pure spinor of $W$ in the left-ideal model, and $\overline N$ that of $\overline W$. Since $I_6$ anticommutes with each $w_k$, $I_6N$ is again annihilated by $W$, hence a multiple of $N$; exactly, $I_6 N = I_7N$ and $I_6\overline N = -I_7\overline N$ (the second is the conjugate of the first, $I_6$ being real). Chirality is the eigenvalue of $I_6$, so $W$ and $\overline W$ are in opposite families, as Proposition \claimref{dark-matter.two-families-are-two-chiralities}(iii) requires of a Witt pair with $m = 3$. Every $J'$ of the same family as $J$ gives the same eigenvalue, by (ii) of that Proposition, so the conclusion holds for every flag with $W\cap\overline W = 0$. In $\Cl(3,3)$, $I_6^2 = (-1)^{15}(+1)^3(-1)^3 = +1$, chirality is real, and conjugation fixes the family; there a flag conjugate cannot be $\kappa$. Checked exactly: isotropy, the Gram matrix, annihilation, the two eigenvalues, $I_6^2$ in both signatures.

\emph{(3) Real-linear exchanges.} Let $g \in O(2,4)$ with $g(W) = \overline W$. Its complex-linear extension is an isometry of $\C^6$ carrying $W$ to a subspace of the other family, so by Proposition \claimref{dark-matter.two-families-are-two-chiralities}(ii) it is the vector action of an odd versor, $\det g = -1$. The integrand $\epsilon(I,V,F,F)$ is the coefficient of $I_6$ in the grade-six part of $IVFF$, a multilinear alternating function of six vector slots, so under the outermorphism of $g$ it is multiplied by $\det g = -1$. An instance is the reflection in $e_1, e_3, e_5$, which carries each $w_k$ to $\overline{w_k}$ and $I_6$ to $-I_6$; on the fixed rational data of the check, $\epsilon(I,V,F,F) = -152$ and its image is $+152$, with or without the composite of grade involution and reversal applied first. The complex-linear exchange $w_k \leftrightarrow v_k$, $v_k = \overline{w_k}/B(w_k,\overline{w_k})$ with $B(w_k,\overline{w_k}) = 2, -2, -2$, sends $e_0 = \tfrac12(w_1 + 2v_1)$ to $\tfrac12 v_1 + w_1$, which has the component $-\tfrac34 I_7e_1$ in $I_7\R^{2,4}$; so in this normalisation it does not preserve $\R^{2,4}$; after the rescaling $w_k \to \tfrac12(1 + I_7)\,w_k$ the dual basis is $v_k = \overline{w_k}, -\overline{w_k}, -\overline{w_k}$ and the exchange is the real map $e_0 \leftrightarrow e_1$, $e_2 \leftrightarrow -e_3$, $e_4 \leftrightarrow -e_5$ of determinant $-1$. In general let $w_k$ be a basis of $W$, $c_{jk} = B(w_j, \overline{w_k})$, and $v_k = \sum_j a_{kj}\overline{w_j}$ the dual basis of $\overline W$, so $a = (c^{T})^{-1}$. The exchange $S$ ($Sw_k = v_k$, $Sv_k = w_k$) sends $\overline{w_k} = \sum_j (a^{-1})_{kj}v_j$ to $\sum_j (a^{-1})_{kj}w_j$, so $\kappa S\kappa\,(w_k) = \sum_j \overline{(a^{-1})_{kj}}\,\overline{w_j}$, and $S$ commutes with $\kappa$, that is preserves $\R^{2,4}$, exactly when $a = \overline{a^{-1}}$, which is $c\,\overline{c} = 1$; for $v_k = \pm\overline{w_k}$, $c$ is diagonal with entries $\pm1$. $S$ is always a complex isometry ($B(Sw_j, Sv_k) = B(v_j, w_k) = \delta_{jk}$, and $W$, $\overline W$ are isotropic) and an involution of trace $0$; when it is real it is a real isometry with $S(W) = \overline W$, so $\det S = -1$ by the first part, and its $-1$ eigenspace is three-dimensional and nondegenerate, so it is a product of three reflections. Checked exactly: the criterion on six bases of this flag (two with $c\,\overline c \neq 1$, four with $c\,\overline c = 1$, one of these with $c$ not diagonal), the image $\tfrac54 e_0 - \tfrac34 I_7e_1$ of $e_0$, and the real map for the rescaled basis.

\emph{(4) The antilinear conjugation.} $\kappa$ is an antilinear automorphism fixing $\Cl(2,4)$ pointwise, and $\mathrm{Ad}(I_6)$ is a linear automorphism with real $I_6$, so they commute and $C = \mathrm{Ad}(I_6)\circ\kappa$ is an antilinear automorphism with $C^2 = \mathrm{Ad}(I_6^2) = \mathrm{Ad}(-1) = \mathrm{id}$. On $\R^{2,4}$ it is $-1$; it fixes $I_6$; it fixes the metric and $\epsilon$ because $-1 \in SO(2,4)$. It carries $N$ to $I_6 \overline N I_6^{-1} = -\overline N$, a pure spinor of $\overline W$. With the flag conjugate realised by $\kappa$, the stored composite is $\widetilde{\ }\circ\alpha\circ\kappa = \widetilde{\ }\circ C$, an antilinear anti-automorphism; on the Fock ideal it sends the left ideal of $W$ to a right ideal of $\overline W$, the dual module, which is the form in which the pairing on the Fock module exchanges grade $k$ with grade $3-k$. A real-linear realisation reverses $I_6$ by (3), and an orientation-preserving one cannot be real-linear by (3), which proves the last sentence of (4). Checked exactly: multiplicativity and antilinearity of $C$ on fixed complex elements, $C^2 = \mathrm{id}$, $C(I_6) = I_6$, $C(N) = -\overline N$.

\emph{(5) The action clause.} In the convention of the formal level, on real tensor fields $\kappa$ acts trivially, so $C$ (and the stored composite, which acts on fields as $C$) acts on every real tensor field of the fibre as the constant element $-1 = \mathrm{Ad}(I_6)$ of $SO_0(2,4)$, by (1); every functional invariant under constant $SO_0(2,4)$ transformations is therefore $C$-invariant. The frozen point's action is $Q(I)^{-1}\int \epsilon(I,V,F,F)$ with $I, V$ fibre vectors and $F$ the curvature, a fibre bivector (\claimref{dark-energy.frozen-point-makes-stelle-west-vector}); under $C$, $I \mapsto -I$, $V \mapsto -V$, $F \mapsto F$ (bivectors are even), $Q(I)$ is unchanged and $\epsilon$ is unchanged, so $S \circ C = S$. With a real-linear flag conjugate $g$ the same computation gives the factor $\det g = -1$ by (3). The base manifold is not acted on: $C$ is a fibre map, and its vector part $-1$ fixes every null line (it sends each null ray to the opposite ray), hence acts trivially on the points of the conformal compactification, the projectivised null cone.

\emph{The convention, and what changes without it.} The stored composite $\phi = \widetilde{\ }\circ\alpha\circ\kappa$ is an anti-automorphism. As a symmetry of fields it acts on vectors by its vector part, $-1$ on real vectors, and on a connection $A$ (bivector valued) by the Lie algebra automorphism $A \mapsto -\phi(A)$ that every anti-automorphism induces on the commutator algebra; on bivectors $-\phi = \alpha = \mathrm{Ad}(I_6)$ is the identity, so $F \mapsto F$, which is the action of $C$. Substituting $\phi$ itself into a bivector slot ($F \mapsto \phi(F) = -F$) is not an action on curvature, since $X \mapsto -X$ is not a Lie algebra automorphism; it agrees with the convention on $\epsilon(I,V,F,F)$, which is even in $F$, but not on functionals odd in a bivector field: on the check's data with $E = \langle ab\rangle_2$ for fixed vectors $a, b$, $\epsilon(I,V,F,E) = -226$ is fixed by $C$ and goes to $+226$ under substitution. A different reading applies $\phi$ to the Clifford integrand $IVFF$ and takes the coefficient of $I_6$: then, since reversal is $-1$ on grade six, the composite with $\kappa$ sends $-152$ to $+152$ and the composite with the real reflection $g$ of (3) keeps $-152$. In that reading the sign statements of (3) to (5) invert and a real-linear flag conjugate composed with Clifford conjugation would preserve the integrand. The statements here are made in the field-action convention, which is how a symmetry acts on fields and how the seventeenth refutation applies $C$ (argumentwise); the computed values of both readings are checked.

\emph{What this does and does not settle.} The premise's description of the conjugation as ``a real linear $\mathbb{Z}_2$ automorphism'' is wrong twice as stated, in the field-action convention: with reversal it is an anti-automorphism, and the flag conjugate compatible with the action clause is antilinear. The antilinear automorphism $A \mapsto \Theta A \Theta^{-1}$ of an antiunitary $\Theta$, and the linear anti-automorphism $A \mapsto \Theta A^{*}\Theta^{-1}$, are the two standard forms of Wigner's antiunitary symmetries (literature, from memory: Wigner 1932; the CPT theorem and Jost's result that total reflection lies in the identity component of the complexified Lorentz group, Streater and Wightman, \emph{PCT, Spin and Statistics, and All That}, 1964, ch.~4); the conformal ambient makes that reflection real by (1). This is offered as the identification, argued, not proved. Clause (5) proves the premise's action clause, in the field-action convention, for functionals of real tensor fields invariant under constant $SO_0(2,4)$. It does not cover spinor fields (real quaternionic spinors of $\Cl(2,4) \cong \Mat_4(\HH)$ are real, but $C$ acts on them by $\mathrm{Ad}(I_6)$, which includes a right multiplication and is not a spin transformation), the internal gauge sector or fermions, on which $\kappa$ acts as charge conjugation on the complex module; there the clause is the reality of the action under $\kappa$, and it stays a premise.
